\question[5]
Ein Medienspieler auf dem Handy hat drei Buttons: Play ($p$), Pause ($a$) und Stop ($s$).
\begin{itemize}
\item Nach dem Start ist der Spieler gestoppt. Mit Play beginnt die Wiedergabe von vorn.
\item Während der Wiedergabe hält Pause an der aktuellen Stelle an. Play setzt die Wiedergabe dort fort.
\item Stop beendet die Wiedergabe und springt zum Anfang zurück, egal ob der Spieler gerade spielt oder pausiert.
\item Ein Button, der im aktuellen Zustand keinen Sinn ergibt (z.B. Pause im gestoppten Zustand), wird ignoriert.
\end{itemize}
\begin{parts}
\part Zeichne das Zustandsdiagramm eines Mealy-Automaten. Die Ausgabe ist die Aktion, die der Spieler ausführt.
\part Gib die Zustandsübergangstabelle an. Ignorierte Eingaben bekommen die Ausgabe --.
\end{parts}

\begin{solutionbox}{8cm}
\begin{parts}
\part Ignorierte Eingaben sind nicht eingezeichnet.

\includegraphics[height=2.5cm]{\pfad/gui_01/gui_01.png}
\part
\begin{tabular}[t]{r|ccc}
 & p & a & s \\ \hline
$\rightarrow$ gestoppt & spielt / Wiedergabe von vorn & gestoppt / -- & gestoppt / -- \\
spielt & spielt / -- & pausiert / anhalten & gestoppt / zum Anfang \\
pausiert & spielt / Wiedergabe fortsetzen & pausiert / -- & gestoppt / zum Anfang \\
\end{tabular}
\end{parts}
\end{solutionbox}
